\begin{textsample}{NEG Dim 6 – global\_north\_2025\_05 – Score -1.00 – global\_north\_2025\_05/125210087.txt}  \label{ex:f6_neg_408}
College of Anaesthesiologists of Ireland faces backlash over guest speaker ' s Gaza-related posts The College of Anaesthesiologists of Ireland ( CAI ) is facing mounting criticism following its decision to invite Professor Idit Matot, an Israeli anesthesiologist, to speak at its 2025 Annual Congress.
The signatories called on the CAI to issue a public statement addressing the ethical implications of the invitation and reaffirming the College ' s commitment to medical neutrality and human rights.
The medical professionals argued that the invitation was inappropriate given the escalating humanitarian crisis in the region and called on the CAI to take action.
It is believed that Professor Matot was scheduled to speak on both Thursday and Friday of the event ( 15th and 16th of May ).
While her first presentation went ahead as planned, controversy erupted shortly afterward when attendees discovered a series of social media posts made by Matot, which were widely perceived as inflammatory and politically charged.
The College of Anaesthesiologists of Ireland.
Pic : Derick P.
Hudson/Shutterstock The discovery prompted immediate calls @ @ @ @ @ @ @ @ @ @ cancelled, to which the CAI complied.
Airing their frustrations, the signatories of the open letter began : ' We, the undersigned, are anaesthesiologists in various stages of their careers are disappointed and appalled at the invitation of Prof.
Idit Matot to the Annual Congress 2025.
Prof.
Matot had posted ' tweets ' during the ongoing Gaza crisis advocating for and attempting to justify the ongoing massacres.
The social media post referenced in the letter sees Hamas stylized with a Nazi swastika alongside a baby with a target on its head holding a teddy bear.
It also includes a map highlighting Israel surrounded by Arab and Muslim-majority countries, implying Israel is under threat and isolated in the region.
The accompanying text claims that European residents could be the next targets of Hamas-style terrorism and warns Europeans not to be complacent or compassionate toward those perceived as threats.
The post also refers to Hamas as ' the new Nazis ' and accuses them of horrific crimes including child murder, rape, and beheading.
It also @ @ @ @ @ @ @ @ @ @ pro-Palestine demonstrators might be terrorists in disguise.
The tweet adds that Palestinians ' must be replaced ' and suggests they can go to ' 22 Arab countries, 57 Islamic countries. ' It emphasizes that Israel is the only Jewish state and asserts it has no interest in Gaza.
It also declares that Israel has already decided how to respond, implying military force.
Further expressing their concerns, the open letter continued : ' We as healthcare workers have a moral and ethical obligation to speak out against injustice and against breaches of international humanitarian law including the targeting of healthcare workers and health facilities and the forced starvation. ' The ongoing crisis in Gaza is condemned by numerous international bodies including the International Criminal Court ( ICC ), the International Court of Justice ( ICJ ), Doctors Without Borders, Amnesty International, and the United Nations.
The ICC and ICJ have highlighted severe breaches of international law by the Israeli army, including war crimes and crimes against humanity.
The United Nations has repeatedly @ @ @ @ @ @ @ @ @ @ colleagues in Gaza, ' the letter continued.
The College of Anaesthesiologists of Ireland.
Pic : Derick P.
Hudson/Shutterstock ' It is imperative that we as anaesthesiologists and healthcare workers advocate for peace no matter where, including the condemnation of terrorism, and the platforming of a person who has publicly expressed such hateful and racist comments is against our moral obligations. ' We request that the College of Anaesthesiologists release an official statement regarding the decision to include this speaker, as well as confirm its commitment to human rights and international law. ' Issuing a response to the signatories this week via email, President Buggy wrote : ' I wish to acknowledge with respect the deeply held views and emotions surrounding the ongoing conflict in the Middle East.
The suffering and apparent slaughter of innocent lives in the region is abhorrent, and I share the longing for peace, justice and dignity for all. ' We strive to be an apolitical, postgraduate medical training body, whose interactions with colleagues are on the @ @ @ @ @ @ @ @ @ @ educational events are invited solely because of their academic and professional expertise in the subject of interest.
Our College and our Annual Scientific Congress are rooted in the values of open inquiry, respectful dialogue and intellectual exchange.
The speaker from Israel, Prof.
Idit Matot, is a recognised clinician scholar and attended in her capacity as an academic and incoming President of our European Society ( ESA-IC ), not as representative of any government or political entity. ' He continued : ' As the congress was beginning on Day 2 on Friday morning, I and the College received multiple complaints about this person contributing to our congress.
Many emanated external to our community but there were also a large number of moderate, sincerely-held protests from across our own community.
Simultaneously, we received credible information indicating that a large-scale, intrusive protest led by outside elements was being planned and therefore, in the interests of everyone ' s safety, I decided to ask Prof.
Matot to withdraw, which she \textbf{kindly} accepted with good grace. ' @ @ @ @ @ @ @ @ @ @ and outrage this invitation caused many of you. ' I abhor violent conflict and am personally angry and appalled at the daily slaughter of innocent lives in Gaza, in Israel in October 2023 and wherever and whenever deaths are caused by violence.
All speaker invitations were given with the intention of supporting open scientific discourse and sharing academic expertise.
However, it is apparent that from now on, professional credentials and open discourse can be only a relative consideration in these matters and College needs to take greater account of contemporary geopolitical circumstances in planning educational events.

% matched lemmas: kindly
\end{textsample}
